%!TeX root = ../main.tex

\subsection{Análisis de controlabilidad del sistema}
La matriz de controlabilidad se define de la forma
\begin{equation}
  \mathcal{C} = \begin{bmatrix}
    \mathbf{B} & \mathbf{A}\mathbf{B} & \cdots & \mathbf{A}^{n-1}\mathbf{B}
  \end{bmatrix} 
\end{equation}

Nuestro sistema corresponde a un orden 3, por lo tanto, su matriz de controlabilidad
está dada por la expresión \eqref{eq:DefiniciónMatrizControlabilidadMotor}

\begin{equation}
  \mathcal{C} = \begin{bmatrix}
    \mathbf{B} & \mathbf{A}\mathbf{B} & \mathbf{A}^{2}\mathbf{B}
  \end{bmatrix} 
  \label{eq:DefiniciónMatrizControlabilidadMotor}
\end{equation}

\begin{equation*}
  \begin{aligned}
    \mathbf{AB} &= 
    \begin{bmatrix}
      -\frac{R}{L}  & 0 & -\frac{K_e}{L} \\
      0             & 0 &  1             \\
      \frac{K_t}{J} & 0 & -\frac{B}{J}   \\
    \end{bmatrix}
    \begin{bmatrix}
      \frac{1}{L} \\
      0           \\
      0           \\
    \end{bmatrix} \\
     &= 
    \begin{bmatrix}
      -\frac{R}{L^2}  \\
      0               \\
      \frac{K_t}{JL}  \\
    \end{bmatrix}
  \end{aligned}
\end{equation*}

\begin{equation*}
  \begin{aligned}
    &\mathbf{A^2B} =\\
    &\begin{bmatrix}
       \tfrac{R^2}{L^2} - \tfrac{K_e K_t}{JL}                 & 0 & \tfrac{K_e}{L}\left(\tfrac{B}{J} + \tfrac{R}{L}\right) \\
       \tfrac{K_t}{J}                                         & 0 & -\tfrac{B}{J}                                          \\
      -\tfrac{K_t}{J}\left(\tfrac{R}{L} + \tfrac{B}{J}\right) & 0 & \tfrac{B^2}{J^2} - \tfrac{K_e K_t}{JL}
    \end{bmatrix}
    \begin{bmatrix}
      \frac{1}{L} \\
      0\\
      0\\
    \end{bmatrix} \\
    &\mathbf{A^2B} =\\
    &\begin{bmatrix}
       \tfrac{R^2}{L^2} - \tfrac{K_e K_t}{JL}                 \\
       \tfrac{K_t}{J}                                         \\
      -\tfrac{K_t}{J}\left(\tfrac{R}{L} + \tfrac{B}{J}\right) \\
    \end{bmatrix}
  \end{aligned}
\end{equation*}

\begin{equation}
  \mathcal{C} = 
  \begin{bmatrix}
    \tfrac{1}{L}     & -\tfrac{R}{L^2}  & \tfrac{R^2}{L^3} - \tfrac{K_e K_t}{JL^2}                  \\
    0                & 0                & \tfrac{K_t}{JL}                                           \\
    0                & \tfrac{K_t}{JL}  & -\tfrac{K_t}{JL}\left( \tfrac{R}{L} + \tfrac{B}{J} \right)
  \end{bmatrix}
  \label{eq:MatrizControlabilidadMotor}
\end{equation}

\subsubsection{Reducción de la matriz de controlabilidad}

\begin{equation*}
  \mathcal{C} = 
  \begin{bmatrix}
      \frac{1}{L} &  -\frac{R}{L^2}   &  \frac{R^2}{L^3}-\frac{K_e K_t}{JL^2}  \\
      0           &   0               &  \frac{K_t}{JL}                        \\
      0           &   \frac{K_t}{JL}  &  -\frac{K_t}{JL}\left(\frac{R}{L}+\frac{B}{J}\right) \\
  \end{bmatrix}
\end{equation*}

\begin{equation*}
  \begin{aligned}
    &R_1\leftarrow L R_1\\
    &\mathcal{C} = 
    \begin{bmatrix}
      1           &  -\frac{R}{L}     &  \frac{R^2}{L^2}-\frac{K_e K_t}{JL}    \\
      0           &   0               &  \frac{K_t}{JL}                        \\
      0           &   \frac{K_t}{JL}  &  -\frac{K_t}{JL}\left(\frac{R}{L}+\frac{B}{J}\right) \\
    \end{bmatrix}
  \end{aligned}
\end{equation*}

\begin{equation*}
  \begin{aligned}
    &R_2 \leftrightarrow R_3 \quad \text{(intercambio de filas)}\\
    &\mathcal{C} = 
    \begin{bmatrix}
      1           &  -\frac{R}{L}     &  \frac{R^2}{L^2}-\frac{K_e K_t}{JL}    \\
      0           &   \frac{K_t}{JL}  &  -\frac{K_t}{JL}\left(\frac{R}{L}+\frac{B}{J}\right) \\
      0           &   0               &  \frac{K_t}{JL}                        \\
    \end{bmatrix}
  \end{aligned}
\end{equation*}

\begin{equation*}
  \begin{aligned}
    &R_2 \leftarrow \tfrac{JL}{K_t} R_2\\
    &\mathcal{C} = 
    \begin{bmatrix}
      1           &  -\frac{R}{L}     &  \frac{R^2}{L^2}-\frac{K_e K_t}{JL}    \\
      0           &   1               &  -\left(\frac{R}{L}+\frac{B}{J}\right) \\
      0           &   0               &  \frac{K_t}{JL}                        \\
    \end{bmatrix}
  \end{aligned}
\end{equation*}

\begin{equation*}
  \begin{aligned}
    &R_3 \leftarrow \tfrac{JL}{K_t} R_3\\
    &\mathcal{C} = 
    \begin{bmatrix}
      1           &  -\frac{R}{L}     &  \frac{R^2}{L^2}-\frac{K_e K_t}{JL}    \\
      0           &   1               &  -\left(\frac{R}{L}+\frac{B}{J}\right) \\
      0           &   0               &  1                                     \\
    \end{bmatrix}
  \end{aligned}
\end{equation*}

\begin{equation*}
  \begin{aligned}
    &R_2 \leftarrow R_2 + \left(\frac{R}{L}+\frac{B}{J}\right) R_3\\
    &\mathcal{C} = 
    \begin{bmatrix}
      1           &  -\frac{R}{L}     &  \frac{R^2}{L^2}-\frac{K_e K_t}{JL}    \\
      0           &   1               &  0                                     \\
      0           &   0               &  1                                     \\
    \end{bmatrix}
  \end{aligned}
\end{equation*}

\begin{equation*}
  \begin{aligned}
    &R_1 \leftarrow R_1 + \tfrac{R}{L} R_2\\
    &\mathcal{C} = 
    \begin{bmatrix}
      1           &   0               &  \frac{R^2}{L^2}-\frac{K_e K_t}{JL}    \\
      0           &   1               &  0                                     \\
      0           &   0               &  1                                     \\
    \end{bmatrix}
  \end{aligned}
\end{equation*}

\begin{equation*}
  \begin{aligned}
    &R_1 \leftarrow R_1 + \left(\tfrac{K_t K_e}{JL}-\tfrac{R^2}{L^2}\right) R_3\\
    &\mathcal{C} = 
    \begin{bmatrix}
      1           &   0               &  0                                     \\
      0           &   1               &  0                                     \\
      0           &   0               &  1                                     \\
    \end{bmatrix}
  \end{aligned}
\end{equation*}

Debido a que el rango de la matriz $\mathcal{C}$ es completo, se puede afirmar que el sistema es controlable.

\subsubsection{Matriz de controlabilidad con valores experimentales}

Con los parámetros experimentales: $R = 2.1\,\Omega$, $L = 0.005\,\text{H}$,
$K_t = 0.12\,\text{N·m/A}$, $K_e = 0.12\,\text{V/(rad/s)}$, $J = 0.0001\,\text{kg·m}^2$
y $B = 1\times10^{-5}\,\text{N·m·s/rad}$, la matriz de controlabilidad numérica es:

\begin{equation*}
  \begin{aligned}
    &\mathcal{C}_{\text{num}} = \\
    &\begin{bmatrix}
      \tfrac{1}{5\mathrm{e}{-3}} & -\tfrac{2.1}{(5\mathrm{e}{-3})^2} & \tfrac{(2.1)^2}{(5\mathrm{e}{-3})^3} - \tfrac{(0.12)^2}{1\mathrm{e}{-4} \cdot (5\mathrm{e}{-3})^2} \\[2pt]
      0 & 0 & \tfrac{0.12}{1\mathrm{e}{-4} \cdot 5\mathrm{e}{-3}} \\[2pt]
      0 & \tfrac{0.12}{1\mathrm{e}{-4} \cdot 5\mathrm{e}{-3}} & -\tfrac{0.12}{1\mathrm{e}{-4} \cdot 5\mathrm{e}{-3}}\left( \tfrac{2.1}{5\mathrm{e}{-3}} + \tfrac{1\mathrm{e}{-5}}{1\mathrm{e}{-4}} \right)
    \end{bmatrix}
  \end{aligned}
\end{equation*}

\begin{equation}
\mathcal{C}_{\text{num}} = 
\begin{bmatrix}
    2.0\mathrm{e}^{2} & -8.4\mathrm{e}^{4} & 2.95\mathrm{e}^{7} \\[2pt]
    0 & 0 & 2.4\mathrm{e}^{5} \\[2pt]
    0 & 2.4\mathrm{e}^{5} & -1.01\mathrm{e}^{8}
\end{bmatrix}
\label{eq:MatrizControlabilidadMotorValoresExperimentales}
\end{equation}
